\section{Contexto general}
Todo coche teledirigido requiere una serie de componentes indispensables, y en este apartado,
se analizarán cada uno de ellos, explicando su función e importancia en el sistema.
Un juguete de estas características lo componen un  \textbf{sistema de comunicación y control},
una \textbf{fuente de energía}, un \textbf{chasis} y un \textbf{sistema de propulsión y dirección}.

\subsection{Comunicación y control}
El sistema de comunicación y control es el encargado de recibir las órdenes del piloto a 
través del mando a distancia y transmitirlas al coche. Este sistema utiliza
tecnologías como radiofrecuencia, Bluetooth o Wi-Fi. Para generar las señales de control
se emplean microcontroladores o microprocesadores, que interpretan las señales recibidas
y las convierten en señales para controlar el motor, la dirección y otros componentes del 
coche. Además, este sistema puede incluir sensores para proporcionar información en tiempo
real sobre la velocidad, posición y situación del coche.

Por lo tanto, se levanta la siguiente pregunta: ¿Qué tecnología de comunicación es la
más adecuada para un coche teledirigido? ¿Cual se va a emplear en este proyecto?
¿Qué tipo de microcontrolador o microprocesador se va a utilizar? Poco a poco se irán
respondiendo a estas preguntas.

\subsubsection{¿Qué es un microcontrolador?}
Un microcontrolador es un dispositivo electrónico integrado que reúne en un único chip
una unidad de procesamiento, memoria y diferentes periféricos de entrada y salida.
Su función principal es ejecutar programas previamente almacenados, permitiendo así
controlar y gestionar el funcionamiento de sistemas electrónicos de forma autónoma.
En el contexto de un coche teledirigido, el microcontrolador es el cerebro del sistema.

\begin{figure}[h]
    \centering
    \includegraphics[width=0.3\textwidth]{../imagenes/Microcontrollers_Atmega32_Atmega8.jpg}
    \caption{Ejemplo de microcontrolador, en este caso, Atmega 32 y Atmega 8}
    \label{fig:microcontroladoratmega}
\end{figure}

Existen diferentes tipos de microcontroladores, cada uno con sus propias características y capacidades:

\paragraph{Arduino.}
Es una plataforma de prototipado electrónico de código 
abierto que utiliza microcontroladores Atmel AVR. Es muy popular 
debido a su facilidad de uso y amplia comunidad de soporte. Se programa
en el lenguaje de programación de Arduino, que es una variante de C/C++ en el
entorno de desarrollo propio de Arduino. Es una buena opción ya que se pueden
añadir módulos de comuncación inalámbrica y el consumo energético es muy
bajo, y es el de menor capacidad de procesamiento de los tres, pero suficiente para 
este proyecto.


\paragraph{ESP32.}
Es un microcontrolador de bajo coste, similar a Arduino, con capacidades de Wi-Fi y 
Bluetooth integradas, lo que lo hace una buen candidato para proyectos y aplicaciones 
que requieren conectividad inalámbrica. El ESP32 es más potente que el Arduino, con 
un procesador de doble núcleo y una mayor cantidad de memoria, lo que permite 
ejecutar tareas más complejas y manejar múltiples sensores y actuadores simultáneamente. 
Se programa utilizando el entorno de desarrollo de Arduino o mediante otros entornos 
compatibles como PlatformIO.

\paragraph{Raspberry Pi:}
Existen dos variantes de Raspberry Pi, la variante microcontrolador (familia Pico)
y la variante microordenador (familia Raspberry Pi). La variante microcontrolador es similar a Arduino,
y se puede programar con microPython o C/C++. Por otro lado, la variante microordenador corre todo un 
sistema operativo Linux por debajo, lo que le permite ejecutar tareas muy complejas como
computación de visión artificial y procesamiento de datos en tiempo real. Es el más potente de todos
pero tambien el más caro y con demasiado consumo energético para este proyecto.


\begin{figure}[H]
    \centering
    
    \begin{subfigure}{0.24\textwidth}
        \centering
        \includegraphics[width=\linewidth]{../imagenes/arduino_nano.jpg}
        \caption{Arduino Nano}
    \end{subfigure}
    \begin{subfigure}{0.24\textwidth}
        \centering
        \includegraphics[width=\linewidth]{../imagenes/ESP32_wroom.jpg}
        \caption{ESP32}
    \end{subfigure}
    \begin{subfigure}{0.24\textwidth}
        \centering
        \includegraphics[width=\linewidth]{../imagenes/Raspberry_Pi_Pico.jpg}
        \caption{Raspberry Pi Pico}
    \end{subfigure}
    \begin{subfigure}{0.22\textwidth}
        \centering
        \includegraphics[width=\linewidth]{../imagenes/Raspberry_Pi_4_Model_B.jpg}
        \caption{Raspberry Pi 4B}
    \end{subfigure}


    \caption{Comparación de microcontroladores}
    \label{fig:comparacion}
\end{figure}


\subsection{Fuente de energía}
Para alimentar un coche teledirigido, se pueden utilizar diferentes tipos de fuentes de
energía, pilas alcalinas, baterías de ion de litio (Li-ion), baterías de polímero de litio (Li-Po)
o baterías de níquel-metal hidruro (NiMH), cada una con sus ventajas y desventajas.

La elección de la fuente de alimentación dependerá de la autonomía deseada, el peso del coche, 
el coste, la facilidad de recarga y la potencia capaz de proporcionar. Las baterías de ion de litio
y polímero de litio son las más comunes en coches teledirigidos debido a su alta densidad energética,
bajo peso y capacidad de recarga rápida. Sin embargo, también son más costosas que las pilas alcalinas
o las baterías de níquel-metal hidruro.


\subsection{Chasis}
El chasis es la estructura que sostiene todos los componentes del coche teledirigido. Puede estar
fabricado con materiales como el plástico, metal, PLA, madera, etc. La estructura debe ser lo
suficientemente resistente como para soportar el peso de los componentes y las fuerzas de aceleración
frenado giros y golpes pero lo suficienteme ligero para no afectar negativamente a la autonomía y 
rendimiento del coche. La distribución del peso de los componentes debe ser equilibrada para que 
el manejo sea óptimo y por último, la estética tambien es importante en un juguete, por lo que debe ser 
atractivo. Por conveniencia y flexibilidad, se optará por un chasis impreso en 3D con PLA, ya que es
un material ligero, resistente y fácil de trabajar, además de permitir diseños personalizados.


\subsection{Propulsión y dirección}
Un coche teledirigido que no se mueve no es un coche teledirigido, ni es nada,
por lo tanto, se necesita un sistema que haga avanzar el coche que estará formado por
un motor de tracción DC, un motor de dirección, un sistema de transmisión y unas ruedas.

\subsubsection{Propulsión:Motor DC}
\subsubsection{Motor de dirección: Servomotor}
\subsubsection{Ruedas}



\cite{Bohren}\cite{Mishchenko}\cite{Bohren,Mishchenko} (prueba, funcionan las citas)
